\section{Introduction}
\label{sec:intro}

Six-player Canadian Fish, also called Literature, is an imperfect-information game played by two
teams of three. Players observe their own cards and the public action history, but not other
players' hands. Because every card transfer is public, a candidate initial deal and the observed
history determine subsequent card ownership. This structure permits exact enumeration of deals
satisfying observer-specific game constraints, providing a reference calculation for
constraint-based inference.

Strategic action choices can convey information beyond legality constraints. We therefore
distinguish the uniform constraint-conditioned distribution from a policy-aware posterior.
\sestina{} uses an approximate belief model with behavioural weighting and combines it with a linear
action policy and guarded search. The study evaluates this configuration against specified internal
baselines, examines component contributions, and assesses sensitivity to partners, rule variants,
and tested adversarial responses.

Computational work on the six-player game appears to be limited. In the literature searches
conducted during this project's development cycles we found no academic work on it, no public
repository containing search, planning or a working learner for it, and no coverage in RLCard
\cite{rlcard} (\S\ref{sec:related-fish}); the scope of those searches is stated there. The nearest
published agent plays a four-player, 48-card variant that does not raise the six-player allocation
problem.

This paper presents \textbf{\sestina{}}, a configuration developed from FishBot \vsix{}, and
evaluates it using the protocol described in \S\ref{sec:evaluation}. Earlier versions established
the engine, the constraint-based inference machinery and a sequence of progressively stronger
policies, but their development also exposed evaluation problems: comparisons that could not resolve
a one-point effect, an optimiser that could not move a policy, and at least two published
conclusions that turned out to be artefacts of the measurement rather than properties of the agent
(\S\ref{sec:background}). The question this paper addresses is whether the tested modifications
improve performance under the revised evaluation protocol.

\paragraph{Version nomenclature.} \sestina{} was developed and evaluated under the identifier
FishBot v0.7. The released agent retains the frozen configuration; source paths, experiment
identifiers, seeds, deal banks, digests and configuration strings retain their original
\texttt{v07} names for reproducibility, including the frozen configuration
\texttt{v07:r12=25,\ldots} of \S\ref{sec:agent} and the sealed banks of \S\ref{sec:evaluation}.
Throughout, \emph{\sestina{}} names the agent and \emph{the v0.7 cycle} names the programme that
produced and measured it. Appendix~\ref{app:repro} gives the detailed mapping.

\subsection{Research questions}
\label{sec:intro-rq}

\begin{description}
\item[RQ1] Does \sestina{} achieve a replicated improvement over \fcheap{}, the registered
  baseline (\S\ref{sec:development-frontier}), measured on material sealed before \sestina{}
  existed?
\item[RQ2] Does that improvement survive changes of partner, cross-play against independently trained
  policies, alternative rule variants, and adversarial responses fitted against it?
\item[RQ3] Which of \sestina{}'s changes account for the improvement, and which candidate
  mechanisms fail to produce measurable strength?
\item[RQ4] What do the results imply about the headroom remaining in the current policy class?
\end{description}

\S\ref{sec:results} answers RQ1 and RQ2 directly, \S\ref{sec:development} and
\S\ref{sec:results-attribution} answer RQ3, and \S\ref{sec:discussion} takes up RQ4.

\subsection{Scope of comparative claims}
\label{sec:intro-strength}

This paper makes comparative claims of three different kinds, and they are kept separate throughout.

\begin{description}
\item[Within-project comparisons.] \sestina{} is compared with its own predecessors and with the
  \emph{\vsix{} frontier}, by which this paper means the set of \vsix{}-era configurations available
  at the end of that cycle: the deployed non-searching policy and the searching configurations
  \fcheap{} and \fmid{}, ordered by search cost (\S\ref{sec:development-frontier}). The paper
  establishes an improvement over the registered baseline \fcheap{}. It does not establish that
  \sestina{} is the strongest configuration in the lineage on every metric: its comparison with a
  development-time composite is unresolved, and on empirical minimax regret over a shared panel it
  ranks third of four (\S\ref{sec:results-panel}).
\item[Robustness evaluation.] Whether the improvement persists against a fixed opponent panel, under
  substituted partners, across rule variants, and against adversarial responses fitted to the
  agent. The paper evaluates these and reports mixed results.
\item[External comparisons.] Whether \sestina{} is the strongest Canadian Fish agent anyone has
  built. The paper does not establish this and does not claim it. The comparison class contains only
  agents this project wrote, and every adversarial-response figure here is the observed edge of a
  particular search rather than a bound on what a search could find.
\end{description}

Accordingly the paper does not describe \sestina{} as near-optimal, unexploitable, or the
strongest known Fish-playing team. \S\ref{sec:discussion-future} states what evidence those claims
would require.

\subsection{Contributions}
\label{sec:intro-contributions}

\begin{enumerate}
\item \textbf{Constraint-based inference} (\S\ref{sec:inference}): a half-suit factorisation and
  count dynamic program for the normalisation constant, card-location marginals, and
  declaration-allocation probabilities under the uniform feasible-deal model.

\item \textbf{Agent specification} (\S\ref{sec:agent}): a reproducible configuration combining
  approximate inference with behavioural weighting, linear action scoring, public-history
  tie-breaking, half-suit contestation, stall detection, and guarded search.

\item \textbf{Evaluation design} (\S\ref{sec:evaluation}): duplicate-deal comparisons, held-out
  deal banks, preregistered criteria, uncertainty estimation, adversarial-search calibration, and
  information-access tests.

\item \textbf{Empirical findings} (\S\ref{sec:results}, \S\ref{sec:development-rejected}): a
  positive comparison with the registered baseline, an unresolved comparison with the development
  composite, component ablations, sensitivity analyses with explicitly stated limitations, and four
  candidate modifications that did not demonstrate an additional improvement under the tested
  conditions.
\end{enumerate}
